% --- Archon physics formalization source begin ---
% archon:physics
% archon:covers PhyXMiniProblems/problem_phyx_mini_0323.lean
% archon:source-report reports/phyx_mini/problem_phyx_mini_0323.source.json

\chapter{Physics problem phyx\_mini\_0323}
\label{ch:PhyXMiniProblems_problem_phyx_mini_0323}

\paragraph{Problem source.}
\#\# Physical scenario

In figure, sound of wavelength \$0.850\$ m is emitted isotropically by point source \$S\$. Sound ray 1 extends directly to detector \$D\$, at distance \$L = 10.0\$ m. Sound ray 2 extends to \$D\$ via a reflection (effectively, a ``bouncing'') of the sound at a flat surface. That reflection occurs on a perpendicular bisector to the \$SD\$ line, at distance \$d\$ from the line. Assume that the reflection shifts the sound wave by \$0.500\textbackslash{}lambda\$.

\#\# Question

For what least value of \$d\$ (other than zero) do the direct sound and the reflected sound arrive at \$D\$ exactly in phase?

\#\# Answer choices

- A: A: 1.18 m

- B: B: 1.23 m

- C: C: 1.35 m

- D: D: 1.47 m

\#\# Figure caption (auxiliary; use the image as primary evidence)

The image depicts a geometrical setup involving light rays and reflective surfaces. Here's a detailed description:

- **Objects and Symbols**:
  - Two circles labeled "S" (Source) and "D" (Detector) at opposite ends.
  - A horizontal line above the circles represents a reflective surface.
  - Two red lines with arrows represent light rays. 
    - Ray 1 travels directly from S to D.
    - Ray 2 reflects off the surface before reaching D.

- **Labels and Text**:
  - "L/2" denotes the horizontal distance from S to the vertical dashed line and from this line to D.
  - "d" is labeled as the vertical height between the reflective surface and the horizontal line containing S and D.
  - The label "Ray 1" is associated with the direct horizontal path, and "Ray 2" with the reflected path.
  - A right angle symbol indicates that the vertical distance "d" is perpendicular to the horizontal line.

- **Relationships**:
  - The setup illustrates how the rays travel from the source to the detector, showing both a direct and a reflected path.
  - The distances are symmetrically divided by the center vertical line, suggesting a scenario like reflection or interference.

This figure is likely used to demonstrate concepts related to optics, such as path difference or interference.

\#\# Dataset metadata

Resonance and Harmonics; Spatial Relation Reasoning, Implicit Condition Reasoning

\paragraph{Recorded answer/context.}
D: D: 1.47 m

\paragraph{Figure/image path.}
/root/proposal\_for\_physic/science-mango/phyx\_mini\_run/phyx\_data/test\_image/323.png

\paragraph{Formalization target.}
create a compiling Lean file with sorry bodies at `PhyXMiniProblems/problem\_phyx\_mini\_0323.lean`. The Lean declarations must preserve the physical quantities, dimensions or dimensional roles, figure labels, governing-law hypotheses, and final relation expressed by this problem.
Use Mathlib/Physlib names found through LeanExplore where available. If a physics API is missing, introduce faithful local abstractions rather than scalar placeholder aliases.

\begin{theorem}[Physics formalization target]
\label{thm:physics:phyx_mini_0323:target}
\lean{PhyXMiniProblems.ProblemPhyXMini0323.problem_phyx_mini_0323}
\uses{def:physics:phyx-mini-0323:phyxminiproblems-problemphyxmini0323-lengthinmeters, def:physics:phyx-mini-0323:phyxminiproblems-problemphyxmini0323-lengthofmeters, def:physics:phyx-mini-0323:phyxminiproblems-problemphyxmini0323-reflectedsoundinterferencesetup, def:physics:phyx-mini-0323:phyxminiproblems-problemphyxmini0323-matchesproblemandfiguredata, def:physics:phyx-mini-0323:phyxminiproblems-problemphyxmini0323-hasphysicalparameters, def:physics:phyx-mini-0323:phyxminiproblems-problemphyxmini0323-satisfiesdepictedraygeometry, def:physics:phyx-mini-0323:phyxminiproblems-problemphyxmini0323-satisfiespropagationandreflectionphaselaw, def:physics:phyx-mini-0323:phyxminiproblems-problemphyxmini0323-isleastpositiveinphaseoffset, def:physics:phyx-mini-0323:phyxminiproblems-problemphyxmini0323-exactleastoffsetinmeters, def:physics:phyx-mini-0323:phyxminiproblems-problemphyxmini0323-recordeddatasetanswer, def:physics:phyx-mini-0323:phyxminiproblems-problemphyxmini0323-roundedtonearesthundredth, def:physics:phyx-mini-0323:phyxminiproblems-problemphyxmini0323-isuniquematchingdisplayedanswer}
The assigned autoformalize agent should translate this physics problem into Lean declarations in the covered file, with theorem and lemma proof bodies written as `by sorry`.
\end{theorem}
\begin{proof}
This is an autoformalization task, not a proof task. Produce statements that can later be proved by the physics prover without weakening the physical model.
\end{proof}

% --- Archon declaration topology begin ---
\section{Declaration topology}
\begin{definition}[Acoustic Length]
  \label{def:physics:phyx-mini-0323:phyxminiproblems-problemphyxmini0323-acousticlength}
  \lean{PhyXMiniProblems.ProblemPhyXMini0323.AcousticLength}
  A signed, unit-independent physical quantity carrying length dimension
\end{definition}
\begin{proof}
  This is the declared model definition of Acoustic Length.
\end{proof}
\begin{definition}[length Readout]
  \label{def:physics:phyx-mini-0323:phyxminiproblems-problemphyxmini0323-lengthreadout}
  \lean{PhyXMiniProblems.ProblemPhyXMini0323.lengthReadout}
  \uses{def:physics:phyx-mini-0323:phyxminiproblems-problemphyxmini0323-acousticlength}
  Read a physical length as a real number in a selected length unit
\end{definition}
\begin{proof}
  This is the declared model definition of length Readout.
\end{proof}
\begin{definition}[length In Meters]
  \label{def:physics:phyx-mini-0323:phyxminiproblems-problemphyxmini0323-lengthinmeters}
  \lean{PhyXMiniProblems.ProblemPhyXMini0323.lengthInMeters}
  \uses{def:physics:phyx-mini-0323:phyxminiproblems-problemphyxmini0323-acousticlength, def:physics:phyx-mini-0323:phyxminiproblems-problemphyxmini0323-lengthreadout}
  The coherent-SI metre readout of an acoustic length
\end{definition}
\begin{proof}
  This is the declared model definition of length In Meters.
\end{proof}
\begin{definition}[length Of Meters]
  \label{def:physics:phyx-mini-0323:phyxminiproblems-problemphyxmini0323-lengthofmeters}
  \lean{PhyXMiniProblems.ProblemPhyXMini0323.lengthOfMeters}
  \uses{def:physics:phyx-mini-0323:phyxminiproblems-problemphyxmini0323-acousticlength}
  Construct a physical length from a prescribed metre readout
\end{definition}
\begin{proof}
  This is the declared model definition of length Of Meters.
\end{proof}
\begin{definition}[Figure Point]
  \label{def:physics:phyx-mini-0323:phyxminiproblems-problemphyxmini0323-figurepoint}
  \lean{PhyXMiniProblems.ProblemPhyXMini0323.FigurePoint}
  The three geometrically distinguished points in the primary figure
\end{definition}
\begin{proof}
  This is the declared model definition of Figure Point.
\end{proof}
\begin{definition}[Sound Ray]
  \label{def:physics:phyx-mini-0323:phyxminiproblems-problemphyxmini0323-soundray}
  \lean{PhyXMiniProblems.ProblemPhyXMini0323.SoundRay}
  The two red acoustic-ray labels printed in the primary figure
\end{definition}
\begin{proof}
  This is the declared model definition of Sound Ray.
\end{proof}
\begin{definition}[Sound Source Kind]
  \label{def:physics:phyx-mini-0323:phyxminiproblems-problemphyxmini0323-soundsourcekind}
  \lean{PhyXMiniProblems.ProblemPhyXMini0323.SoundSourceKind}
  Physical kind of source specified in the scenario
\end{definition}
\begin{proof}
  This is the declared model definition of Sound Source Kind.
\end{proof}
\begin{definition}[Receiver Kind]
  \label{def:physics:phyx-mini-0323:phyxminiproblems-problemphyxmini0323-receiverkind}
  \lean{PhyXMiniProblems.ProblemPhyXMini0323.ReceiverKind}
  Physical kind of receiving object specified in the scenario
\end{definition}
\begin{proof}
  This is the declared model definition of Receiver Kind.
\end{proof}
\begin{definition}[Reflector Kind]
  \label{def:physics:phyx-mini-0323:phyxminiproblems-problemphyxmini0323-reflectorkind}
  \lean{PhyXMiniProblems.ProblemPhyXMini0323.ReflectorKind}
  Physical kind of boundary at which ray 2 bounces
\end{definition}
\begin{proof}
  This is the declared model definition of Reflector Kind.
\end{proof}
\begin{definition}[Figure Feature]
  \label{def:physics:phyx-mini-0323:phyxminiproblems-problemphyxmini0323-figurefeature}
  \lean{PhyXMiniProblems.ProblemPhyXMini0323.FigureFeature}
  Qualitative labels and relations visible in the primary figure
\end{definition}
\begin{proof}
  This is the declared model definition of Figure Feature.
\end{proof}
\begin{definition}[ray Route]
  \label{def:physics:phyx-mini-0323:phyxminiproblems-problemphyxmini0323-rayroute}
  \lean{PhyXMiniProblems.ProblemPhyXMini0323.rayRoute}
  \uses{def:physics:phyx-mini-0323:phyxminiproblems-problemphyxmini0323-figurepoint, def:physics:phyx-mini-0323:phyxminiproblems-problemphyxmini0323-soundray}
  Ordered route through the labelled points followed by each sound ray
\end{definition}
\begin{proof}
  This is the declared model definition of ray Route.
\end{proof}
\begin{definition}[Reflected Sound Interference Setup]
  \label{def:physics:phyx-mini-0323:phyxminiproblems-problemphyxmini0323-reflectedsoundinterferencesetup}
  \lean{PhyXMiniProblems.ProblemPhyXMini0323.ReflectedSoundInterferenceSetup}
  \uses{def:physics:phyx-mini-0323:phyxminiproblems-problemphyxmini0323-acousticlength, def:physics:phyx-mini-0323:phyxminiproblems-problemphyxmini0323-soundray, def:physics:phyx-mini-0323:phyxminiproblems-problemphyxmini0323-soundsourcekind, def:physics:phyx-mini-0323:phyxminiproblems-problemphyxmini0323-receiverkind, def:physics:phyx-mini-0323:phyxminiproblems-problemphyxmini0323-reflectorkind, def:physics:phyx-mini-0323:phyxminiproblems-problemphyxmini0323-figurefeature}
  The independent physical data and observable behavior of the interference experiment. The argument of rayPathLengthInMetersAtOffset is a physical candidate offset d. arriveExactlyInPhaseAtOffset remains an abstract acoustic observation here and is related to path and reflection phase by the governing-law structure below. No field assigns the requested least positive value of d
\end{definition}
\begin{proof}
  This is the declared model definition of Reflected Sound Interference Setup.
\end{proof}
\begin{definition}[Matches Problem And Figure Data]
  \label{def:physics:phyx-mini-0323:phyxminiproblems-problemphyxmini0323-matchesproblemandfiguredata}
  \lean{PhyXMiniProblems.ProblemPhyXMini0323.MatchesProblemAndFigureData}
  \uses{def:physics:phyx-mini-0323:phyxminiproblems-problemphyxmini0323-lengthinmeters, def:physics:phyx-mini-0323:phyxminiproblems-problemphyxmini0323-reflectedsoundinterferencesetup}
  Scenario and primary-image data. The wavelength is 0.850 m = 17/20 m, the source-detector distance is 10.0 m, and reflection shifts the wave by 0.500 lambda, equivalently one half-cycle. This predicate does not select any value of the unknown offset d
\end{definition}
\begin{proof}
  This is the declared model definition of Matches Problem And Figure Data.
\end{proof}
\begin{definition}[Has Physical Parameters]
  \label{def:physics:phyx-mini-0323:phyxminiproblems-problemphyxmini0323-hasphysicalparameters}
  \lean{PhyXMiniProblems.ProblemPhyXMini0323.HasPhysicalParameters}
  \uses{def:physics:phyx-mini-0323:phyxminiproblems-problemphyxmini0323-lengthinmeters, def:physics:phyx-mini-0323:phyxminiproblems-problemphyxmini0323-reflectedsoundinterferencesetup}
  Positivity assumptions for the independently supplied physical lengths
\end{definition}
\begin{proof}
  This is the declared model definition of Has Physical Parameters.
\end{proof}
\begin{definition}[Satisfies Depicted Ray Geometry]
  \label{def:physics:phyx-mini-0323:phyxminiproblems-problemphyxmini0323-satisfiesdepictedraygeometry}
  \lean{PhyXMiniProblems.ProblemPhyXMini0323.SatisfiesDepictedRayGeometry}
  \uses{def:physics:phyx-mini-0323:phyxminiproblems-problemphyxmini0323-acousticlength, def:physics:phyx-mini-0323:phyxminiproblems-problemphyxmini0323-lengthinmeters, def:physics:phyx-mini-0323:phyxminiproblems-problemphyxmini0323-reflectedsoundinterferencesetup}
  The two path-length laws read from the symmetric figure. Ray 1 has length L. Ray 2 consists of two congruent right-triangle hypotenuses with legs L/2 and d, giving 2 * sqrt ((L/2)\textasciicircum{}2 + d\textasciicircum{}2). These equations apply to every nonnegative proposed offset and do not state which offset yields the first in-phase arrival
\end{definition}
\begin{proof}
  This is the declared model definition of Satisfies Depicted Ray Geometry.
\end{proof}
\begin{definition}[Satisfies Propagation And Reflection Phase Law]
  \label{def:physics:phyx-mini-0323:phyxminiproblems-problemphyxmini0323-satisfiespropagationandreflectionphaselaw}
  \lean{PhyXMiniProblems.ProblemPhyXMini0323.SatisfiesPropagationAndReflectionPhaseLaw}
  \uses{def:physics:phyx-mini-0323:phyxminiproblems-problemphyxmini0323-acousticlength, def:physics:phyx-mini-0323:phyxminiproblems-problemphyxmini0323-lengthinmeters, def:physics:phyx-mini-0323:phyxminiproblems-problemphyxmini0323-reflectedsoundinterferencesetup}
  For monochromatic propagation, path excess divided by wavelength is a phase difference measured in cycles. Reflection adds its half-cycle shift. The arrivals are exactly in phase precisely when the total cycle difference is an integer. This is a general governing law for every nonnegative offset, not a premise asserting the requested answer
\end{definition}
\begin{proof}
  This is the declared model definition of Satisfies Propagation And Reflection Phase Law.
\end{proof}
\begin{definition}[Is Least Positive In Phase Offset]
  \label{def:physics:phyx-mini-0323:phyxminiproblems-problemphyxmini0323-isleastpositiveinphaseoffset}
  \lean{PhyXMiniProblems.ProblemPhyXMini0323.IsLeastPositiveInPhaseOffset}
  \uses{def:physics:phyx-mini-0323:phyxminiproblems-problemphyxmini0323-acousticlength, def:physics:phyx-mini-0323:phyxminiproblems-problemphyxmini0323-lengthinmeters, def:physics:phyx-mini-0323:phyxminiproblems-problemphyxmini0323-reflectedsoundinterferencesetup}
  An offset answers the question when it is strictly positive, produces exactly in-phase arrivals, and is no larger than any other strictly positive offset with that property. The explicit minimality clause formalizes "least value other than zero"
\end{definition}
\begin{proof}
  This is the declared model definition of Is Least Positive In Phase Offset.
\end{proof}
\begin{definition}[exact Least Offset In Meters]
  \label{def:physics:phyx-mini-0323:phyxminiproblems-problemphyxmini0323-exactleastoffsetinmeters}
  \lean{PhyXMiniProblems.ProblemPhyXMini0323.exactLeastOffsetInMeters}
  \uses{def:physics:phyx-mini-0323:phyxminiproblems-problemphyxmini0323-lengthinmeters, def:physics:phyx-mini-0323:phyxminiproblems-problemphyxmini0323-reflectedsoundinterferencesetup}
  Solving the first constructive condition, whose path excess is lambda/2, gives the exact metre expression d = sqrt (lambda * (4 * L + lambda)) / 4. This definition only names the candidate expression; proving that it is the least positive in-phase offset remains the substantive conclusion below
\end{definition}
\begin{proof}
  This is the declared model definition of exact Least Offset In Meters.
\end{proof}
\begin{definition}[Answer Choice]
  \label{def:physics:phyx-mini-0323:phyxminiproblems-problemphyxmini0323-answerchoice}
  \lean{PhyXMiniProblems.ProblemPhyXMini0323.AnswerChoice}
  Labels of the four displayed answer choices
\end{definition}
\begin{proof}
  This is the declared model definition of Answer Choice.
\end{proof}
\begin{definition}[displayed Offset In Meters]
  \label{def:physics:phyx-mini-0323:phyxminiproblems-problemphyxmini0323-displayedoffsetinmeters}
  \lean{PhyXMiniProblems.ProblemPhyXMini0323.displayedOffsetInMeters}
  \uses{def:physics:phyx-mini-0323:phyxminiproblems-problemphyxmini0323-answerchoice}
  Metre value printed beside each answer label
\end{definition}
\begin{proof}
  This is the declared model definition of displayed Offset In Meters.
\end{proof}
\begin{definition}[recorded Dataset Answer]
  \label{def:physics:phyx-mini-0323:phyxminiproblems-problemphyxmini0323-recordeddatasetanswer}
  \lean{PhyXMiniProblems.ProblemPhyXMini0323.recordedDatasetAnswer}
  \uses{def:physics:phyx-mini-0323:phyxminiproblems-problemphyxmini0323-answerchoice}
  The answer label recorded in the supplied dataset metadata
\end{definition}
\begin{proof}
  This is the declared model definition of recorded Dataset Answer.
\end{proof}
\begin{definition}[rounded To Nearest Hundredth]
  \label{def:physics:phyx-mini-0323:phyxminiproblems-problemphyxmini0323-roundedtonearesthundredth}
  \lean{PhyXMiniProblems.ProblemPhyXMini0323.roundedToNearestHundredth}
  Round a metre readout to the nearest hundredth of a metre
\end{definition}
\begin{proof}
  This is the declared model definition of rounded To Nearest Hundredth.
\end{proof}
\begin{definition}[Matches Displayed Answer]
  \label{def:physics:phyx-mini-0323:phyxminiproblems-problemphyxmini0323-matchesdisplayedanswer}
  \lean{PhyXMiniProblems.ProblemPhyXMini0323.MatchesDisplayedAnswer}
  \uses{def:physics:phyx-mini-0323:phyxminiproblems-problemphyxmini0323-reflectedsoundinterferencesetup, def:physics:phyx-mini-0323:phyxminiproblems-problemphyxmini0323-exactleastoffsetinmeters, def:physics:phyx-mini-0323:phyxminiproblems-problemphyxmini0323-answerchoice, def:physics:phyx-mini-0323:phyxminiproblems-problemphyxmini0323-displayedoffsetinmeters, def:physics:phyx-mini-0323:phyxminiproblems-problemphyxmini0323-roundedtonearesthundredth}
  A choice matches the least positive offset at the displayed precision
\end{definition}
\begin{proof}
  This is the declared model definition of Matches Displayed Answer.
\end{proof}
\begin{definition}[Is Unique Matching Displayed Answer]
  \label{def:physics:phyx-mini-0323:phyxminiproblems-problemphyxmini0323-isuniquematchingdisplayedanswer}
  \lean{PhyXMiniProblems.ProblemPhyXMini0323.IsUniqueMatchingDisplayedAnswer}
  \uses{def:physics:phyx-mini-0323:phyxminiproblems-problemphyxmini0323-reflectedsoundinterferencesetup, def:physics:phyx-mini-0323:phyxminiproblems-problemphyxmini0323-answerchoice, def:physics:phyx-mini-0323:phyxminiproblems-problemphyxmini0323-matchesdisplayedanswer}
  The selected label is the unique displayed choice matching the result
\end{definition}
\begin{proof}
  This is the declared model definition of Is Unique Matching Displayed Answer.
\end{proof}
\begin{lemma}[exact offset is least positive in phase]
  \label{lem:physics:phyx-mini-0323:phyxminiproblems-problemphyxmini0323-exact-offset-is-least-positive-in-phase}
  \lean{PhyXMiniProblems.ProblemPhyXMini0323.exact_offset_is_least_positive_in_phase}
  \uses{def:physics:phyx-mini-0323:phyxminiproblems-problemphyxmini0323-lengthofmeters, def:physics:phyx-mini-0323:phyxminiproblems-problemphyxmini0323-reflectedsoundinterferencesetup, def:physics:phyx-mini-0323:phyxminiproblems-problemphyxmini0323-matchesproblemandfiguredata, def:physics:phyx-mini-0323:phyxminiproblems-problemphyxmini0323-hasphysicalparameters, def:physics:phyx-mini-0323:phyxminiproblems-problemphyxmini0323-satisfiesdepictedraygeometry, def:physics:phyx-mini-0323:phyxminiproblems-problemphyxmini0323-satisfiespropagationandreflectionphaselaw, def:physics:phyx-mini-0323:phyxminiproblems-problemphyxmini0323-isleastpositiveinphaseoffset, def:physics:phyx-mini-0323:phyxminiproblems-problemphyxmini0323-exactleastoffsetinmeters}
  The ray geometry and phase law imply that the exact square-root expression is the least positive offset for which the two arrivals are in phase
\end{lemma}
\begin{proof}
  The conclusion follows from the governing relations and figure readouts named in the statement of exact offset is least positive in phase.
\end{proof}
% --- Archon declaration topology end ---
% --- Archon physics formalization source end ---
